\scp{Issues in methodology}{02-05}
An assumption of a monolithic spread of the Austronesians into
the region can easily make it difficult to see the discontinuities
that are evident in the data.
An assumption of a single higher level subgroup encompassing
most or all of the Austronesian languages of our target region
can make it hard to see that sporadically similar patterns are
not actually exclusively shared innovations,
once the sound correspondences are mapped out
and fuller data are examined.

Conversely, an assumption that there is not a single encompassing
subgroup in a region can make it difficult to see similarities
and links that are actually there.
An assumption that lower level subgrouping arguments posited
by others “must be correct”,
needs to be held with caution.
We ourselves found several cases in which,
out of respect for the work of others,
we initially adopted their subgroups
and worked with those.
But in following where the data led us,
(often with additional, more complete
and more reliable data),
we found there were other issues that had to be given greater
weight, and eventually we had to depart from earlier classifications,
redefine the Wallacean subgroups (or nodes within them),
and even rewrite and reorganise chapters already drafted in this book.

We also need to be cautious about jumping to conclusions about
time depth based on the linguistic evidence.
While certain sound changes can inform us about 
a relative sequence of which ones 
came first and which ones came later
(e.g.
*b > \it{f > h},
*p > \it{ɸ > h >} {\0},
*k > \it{ʔ} > {\0},
*s > \it{h}) on their own they cannot tell us at
what point in time these changes happened.
%While certain sound changes can tell us about a relative sequence
%of what most likely was first,
%and what came later (e.g.
%*b > \it{f > h},
%*p > \it{ɸ > h >} {\0},
%*k > \it{ʔ} > {\0},
%*s > \it{h}), on their own they cannot tell us when that happened.
The field of glottochronology has been largely abandoned because
we know, for example, that a) languages do not change at a constant rate
through time, b) different languages change at different rates,
and c) there can be spurts of rapid change,
for example, associated with war
and conquest or tragedy such in the aftermath of volcanic eruptions and tsunamis.
But sometimes a datable change is extrapolated to other languages,
and this may or may not be a problem.
For example, earlier it was noted that \citet[10]{Klamer2017} linked
the break-up of the \tsc{\ili{Alor}-Pantar} Papuan languages to around
the same time as the arrival of the Austronesians because “the
reconstructed vocabulary of proto-AP appears to contain Austronesian
loan words such as ‘betel nut’ [{\ldots}] following regular sound
changes suggest that the AP family split up after being in contact
with the Austronesian languages in the area.”

That is certainly possible,
but \citet[81]{Grimes1991c} has also noted that several Portuguese
loans in \ili{Buru} display the same changes to sounds
and word structure that are found in PMP words.
So Portuguese \it{meirinho} /maiɾiɲu/ ‘bailiff’ > \ili{Ambon} \it{marinyu}
> \ili{Buru} \it{emrimo} ‘village crier, constable’.
And Portuguese \it{capitão} /kapitãũ/ ‘ship captain,
military officer’ > \ili{Ambon} \it{kapitan} > \ili{Buru} \it{epkitan} ‘fighting champion, wartime leader’.
Similar restructuring is seen in \ili{Buru} loans ultimately from \ili{Sanskrit},
such as: \it{baksa} ‘read’,
\it{eʔwasa-t} ‘wealth, power, influence’, \it{fisara} ‘speak’
(\ili{Lisela} \it{epsara}), \it{mansari} ‘hunt’, \it{mansia} ‘human’,
and \it{storita} ‘story, tell a story’.
And also in \ili{Buru} loans ultimately from \ili{Arabic},
such as: \it{barkate} ‘bless (v),
blessing (n)’, \it{eʔʤadi} ‘become,
be born’, \it{eʔrama-t} ‘devout,
spiritual’, \it{eslau} ‘beg forgiveness, restore right relationships’,
\it{haiwane} ‘domestic animals’, \it{hasa-k} ‘oppose,
hate’, \it{hum sikit} ‘mosque’, \it{ʤunai} ‘world’,
\it{kufur} ‘grave’, \it{slame} ‘Islam,
Muslim’, \it{srane} ‘Christian, christen’.

\citet[397]{Edwards2021} could reconstruct a Proto-Rote-\ili{Meto} form **ŋgarei ‘church’
based on regular sound correspondences.
Yet he has not done so,
since he is aware that the daughter forms are from Portuguese \it{igreja} [iˈgreʒa].
Portuguese contact did not happen before 1512 AD,
and \ili{Sanskrit} and \ili{Arabic} influence came significantly later than
Austronesian incursion, while the connection to PMP in the region goes back about 4,000 years.
So parallel developments or processes in the phonotactics
and historical phonology may not necessarily reflect similarities in time depth.

Papuan-sourced {\#}muku ‘banana’ and Austronesian-sourced reconstructions
such as ‘betelnut’ in Papuan TAP languages provide evidence that
words can and do spread across unrelated groups that are not sourced in
a common parent language or linkage.
Additional similar examples are given in the later chapters of this book.
In a geographically diverse
and dispersed region with known trade
and travel, these examples should encourage serious caution (and
scepticism) for using lexical evidence as the basis for subgrouping arguments.

